\section{สรุปผลการดำเนินงาน}
\sloppy

บทนี้จะกล่าวถึงบทสรุปของการดำเนินงานโครงงานระบบบริหารจัดการยิม (GymBro Management System) โดยครอบคลุมถึงภาพรวมของระบบที่พัฒนาขึ้น ผลลัพธ์ที่ได้จากการพัฒนา ปัญหาและอุปสรรคที่พบระหว่างการดำเนินงาน รวมถึงข้อเสนอแนะและแนวทางในการพัฒนาต่อยอดระบบในอนาคต

\subsection{สรุปภาพรวมของระบบ (System Overview Summary)}
ระบบบริหารจัดการยิม (GymBro Management System) ถูกพัฒนาขึ้นเพื่อเป็นเว็บแอปพลิเคชันสำหรับช่วยบริหารจัดการข้อมูลและการดำเนินงานภายในยิมให้มีความสะดวก รวดเร็ว และถูกต้องแม่นยำมากยิ่งขึ้น โดยมีวัตถุประสงค์หลักเพื่อลดความผิดพลาดจากการทำงานด้วยมือ (Manual Process) และเพิ่มประสิทธิภาพในการให้บริการแก่สมาชิก

ระบบถูกออกแบบให้รองรับการทำงานของกลุ่มผู้ใช้งาน 2 กลุ่มหลัก ได้แก่
\begin{enumerate}
    \item \textbf{ผู้ดูแลระบบ (Administrator):} สามารถบริหารจัดการข้อมูลพื้นฐานของระบบ เช่น ข้อมูลสมาชิก (Customers), ข้อมูลเทรนเนอร์ (Trainers), ข้อมูลอุปกรณ์ออกกำลังกาย (Gym Equipments) และการจัดการตารางการจองคลาสเรียน (Training Sessions) รวมถึงสามารถดูบันทึกการใช้งานระบบ (System Logs) ได้
    \item \textbf{สมาชิกทั่วไป (Member/User):} สามารถเข้าสู่ระบบเพื่อดูข้อมูลส่วนตัว แก้ไขโปรไฟล์ และตรวจสอบสถานะการเป็นสมาชิกได้
\end{enumerate}

ในส่วนของสถาปัตยกรรมระบบ ได้มีการเลือกใช้เทคโนโลยี Full Stack Development ที่ทันสมัย โดยฝั่ง Server-Side ใช้ Node.js ร่วมกับ Express.js ในการสร้าง RESTful API และ Web Server ส่วนฝั่ง Client-Side ใช้ EJS (Embedded JavaScript Templating) ร่วมกับ Tailwind CSS ในการแสดงผลหน้าเว็บที่สวยงามและรองรับการใช้งานบนอุปกรณ์ต่างๆ (Responsive Design) และใช้ฐานข้อมูล PostgreSQL ในการจัดเก็บข้อมูลที่มีความสัมพันธ์กัน (Relational Database) ผ่าน Sequelize ORM

\subsection{ผลการพัฒนาระบบ (Development Achievements)}
จากการดำเนินงานพัฒนาโครงงาน สามารถสรุปผลการพัฒนาระบบในส่วนต่างๆ ได้ดังนี้:

\begin{itemize}
    \item \textbf{ระบบจัดการฐานข้อมูล (Database Management):} สามารถออกแบบและสร้างฐานข้อมูลเชิงสัมพันธ์ (Relational Database) ที่มีความถูกต้องตามหลักการ Normalization และสามารถเชื่อมต่อกับแอปพลิเคชันผ่าน Sequelize ORM ได้อย่างสมบูรณ์ รองรับการทำ CRUD Operations (Create, Read, Update, Delete) กับข้อมูลทุกตาราง
    
    \item \textbf{ระบบ Backend API:} ได้พัฒนา RESTful API ที่มีประสิทธิภาพสำหรับการจัดการข้อมูลหลักของระบบ ได้แก่:
    \begin{itemize}
        \item API สำหรับการเข้าสู่ระบบ (Authentication) และการตรวจสอบสิทธิ์ (Authorization)
        \item API สำหรับจัดการข้อมูลสมาชิก (Customers), เทรนเนอร์ (Trainers), และอุปกรณ์ (Equipments)
        \item API สำหรับการจองคลาสเรียน (Training Sessions) และการคำนวณวันหมดอายุสมาชิก
    \end{itemize}

    \item \textbf{ระบบ Frontend Interface:} ได้พัฒนาส่วนติดต่อผู้ใช้งาน (User Interface) ที่มีความสวยงามและใช้งานง่าย โดยแบ่งส่วนการทำงานระหว่างหน้าบ้าน (User Dashboard) และหลังบ้าน (Admin Dashboard) อย่างชัดเจน มีการใช้ JavaScript ในการดึงข้อมูลจาก API มาแสดงผลแบบ Dynamic และมีการจัดการ Form Validation เพื่อป้องกันข้อมูลที่ผิดพลาด

    \item \textbf{ระบบความปลอดภัยเบื้องต้น:} มีการตรวจสอบสิทธิ์การเข้าถึงข้อมูล (Role-Based Access Control) โดยแยกสิทธิ์ระหว่าง Admin และ User อย่างชัดเจน และมีการป้องกันการเข้าถึงหน้าเว็บที่ไม่ได้รับอนุญาต (Route Guarding)
    
    \item \textbf{ฟังก์ชันพิเศษ:} ระบบมีการบันทึกประวัติการใช้งาน \newline (Logging) ลงในไฟล์ JSON เพื่อให้ผู้ดูแลระบบสามารถตรวจสอบย้อนหลังได้ และมีฟังก์ชันการเรียงลำดับ ID ของข้อมูลใหม่ \newline (ID Reordering) เพื่อให้ข้อมูลมีความต่อเนื่องสวยงามเมื่อมีการลบข้อมูลออกไป
\end{itemize}

\subsection{ปัญหาและข้อจำกัด (Limitations and Challenges)}
จากการวิเคราะห์และทดสอบระบบ พบว่ายังมีข้อจำกัดและประเด็นที่ควรได้รับการปรับปรุง ดังนี้:

\begin{enumerate}
    \item \textbf{ความปลอดภัยของข้อมูล (Security):} \newline
    \begin{itemize}
        \item ระบบการเข้าสู่ระบบปัจจุบันยังใช้การตรวจสอบข้อมูลพื้นฐาน (Email และ Phone Number) ซึ่งอาจไม่เพียงพอสำหรับระบบที่มีความสำคัญสูง ควรมีการเข้ารหัสรหัสผ่าน (Password Hashing) และใช้ Token-based Authentication (เช่น JWT) เพื่อความปลอดภัยที่มากขึ้น
        \item ข้อมูลสิทธิ์ผู้ดูแลระบบ (Admin Credentials) บางส่วนยังมีการกำหนดค่าคงที่ (Hardcoded) ในซอร์สโค้ด ซึ่งควรย้ายไปเก็บใน Environment Variables หรือฐานข้อมูล
    \end{itemize}

    \item \textbf{ประสิทธิภาพและการขยายตัว (Scalability):} \newline
    \begin{itemize}
        \item การบันทึก Logs ลงในไฟล์ JSON (File-based Logging) อาจทำให้เกิดปัญหาด้านประสิทธิภาพเมื่อไฟล์มีขนาดใหญ่ขึ้น ควรเปลี่ยนไปใช้การบันทึกลงฐานข้อมูลหรือบริการ Logging ภายนอก
        \item การเรียงลำดับ ID ใหม่ (ID Reordering) ทุกครั้งที่มีการลบข้อมูล เป็นกระบวนการที่ใช้ทรัพยากรสูงและอาจเกิดปัญหา Data Integrity หากมีผู้ใช้งานจำนวนมากพร้อมกัน
    \end{itemize}

    \item \textbf{ฟีเจอร์การทำงาน (Feature Completeness):} \newline
    \begin{itemize}
        \item ระบบยังขาดฟังก์ชันการแจ้งเตือนแบบ Real-time (เช่น การแจ้งเตือนเมื่อการจองสำเร็จ)
        \item ยังไม่มีระบบการชำระเงิน (Payment Gateway) สำหรับการสมัครสมาชิกหรือจองคลาสแบบเสียค่าใช้จ่าย
    \end{itemize}
\end{enumerate}

\subsection{แนวทางในการพัฒนาต่อยอด (Future Work)}
เพื่อให้ระบบมีความสมบูรณ์และรองรับการใช้งานจริงได้ดียิ่งขึ้น ผู้จัดทำขอเสนอแนวทางในการพัฒนาต่อยอดในอนาคต ดังนี้:

\begin{enumerate}
    \item \textbf{ยกระดับมาตรการความปลอดภัย:} \newline พัฒนาระบบยืนยันตัวตนให้รัดกุมยิ่งขึ้น โดยใช้มาตรฐาน OAuth 2.0 หรือ JWT และเข้ารหัสข้อมูลสำคัญในฐานข้อมูล รวมถึงการทำ HTTPS เพื่อความปลอดภัยในการรับส่งข้อมูล
    
    \item \textbf{เพิ่มระบบการจัดการทางการเงิน:} \newline เชื่อมต่อกับ Payment Gateway \newline (เช่น Stripe หรือ Omise) เพื่อรองรับการชำระค่าสมาชิกและค่าบริการเทรนเนอร์ผ่านบัตรเครดิตหรือ QR Code
    
    \item \textbf{พัฒนา Mobile Application:} \newline สร้างแอปพลิเคชันบนมือถือ (iOS/Android) เพื่อให้สมาชิกสามารถจองคลาสและตรวจสอบตารางเรียนได้สะดวกยิ่งขึ้น โดยใช้ API เดิมที่มีอยู่
    
    \item \textbf{ปรับปรุงระบบ Dashboard:} \newline เพิ่มการแสดงผลข้อมูลเชิงลึก (Data Analytics) เช่น กราฟแสดงจำนวนผู้เข้าใช้บริการในแต่ละช่วงเวลา หรือยอดรายได้ประจำเดือน เพื่อช่วยในการตัดสินใจของผู้บริหาร
    
    \item \textbf{ระบบการแจ้งเตือน (Notification System):} \newline เพิ่มระบบแจ้งเตือนผ่าน Email หรือ LINE Notify เพื่อแจ้งเตือนสมาชิกเมื่อใกล้ถึงเวลาจองคลาส หรือเมื่อสถานะสมาชิกใกล้หมดอายุ
\end{enumerate}
